\documentclass[10pt,a4paper]{amsart}
\usepackage{amsmath,amssymb,amsthm,graphicx}
\usepackage{verbatim}
\newenvironment{reference}{\comment}{\endcomment}
%\newenvironment{reference}{}{}
\newenvironment{slogan}{\comment}{\endcomment}
\newenvironment{history}{\comment}{\endcomment}

% For commutative diagrams we use Xy-pic
\usepackage[all]{xy}

% We use 2cell for 2-commutative diagrams.
\xyoption{2cell}
% Theorem environments.
%
\theoremstyle{plain}
\newtheorem{theorem}[subsection]{Theorem}
\newtheorem{proposition}[subsection]{Proposition}
\newtheorem{lemma}[subsection]{Lemma}

\theoremstyle{definition}
\newtheorem{definition}[subsection]{Definition}
\newtheorem{example}[subsection]{Example}
\newtheorem{exercise}[subsection]{Exercise}
\newtheorem{situation}[subsection]{Situation}

\theoremstyle{remark}
\newtheorem{remark}[subsection]{Remark}
\newtheorem{remarks}[subsection]{Remarks}

\numberwithin{equation}{subsection}

% Macros
%
\def\lim{\mathop{\mathrm{lim}}\nolimits}
\def\colim{\mathop{\mathrm{colim}}\nolimits}
\def\Spec{\mathop{\mathrm{Spec}}}
\def\Hom{\mathop{\mathrm{Hom}}\nolimits}
\def\Ext{\mathop{\mathrm{Ext}}\nolimits}
\def\SheafHom{\mathop{\mathcal{H}\!\mathit{om}}\nolimits}
\def\SheafExt{\mathop{\mathcal{E}\!\mathit{xt}}\nolimits}
\def\Sch{\mathit{Sch}}
\def\Mor{\mathop{\mathrm{Mor}}\nolimits}
\def\Ob{\mathop{\mathrm{Ob}}\nolimits}
\def\Sh{\mathop{\mathit{Sh}}\nolimits}
\def\NL{\mathop{N\!L}\nolimits}
\def\CH{\mathop{\mathrm{CH}}\nolimits}
\def\proetale{{pro\text{-}\acute{e}tale}}
\def\etale{{\acute{e}tale}}
\def\QCoh{\mathit{QCoh}}
\def\Ker{\mathop{\mathrm{Ker}}}
\def\Im{\mathop{\mathrm{Im}}}
\def\Coker{\mathop{\mathrm{Coker}}}
\def\Coim{\mathop{\mathrm{Coim}}}

% Boxtimes
%
\DeclareMathSymbol{\boxtimes}{\mathbin}{AMSa}{"02}

%
% Macros for moduli stacks/spaces
%
\def\QCohstack{\mathcal{QC}\!\mathit{oh}}
\def\Cohstack{\mathcal{C}\!\mathit{oh}}
\def\Spacesstack{\mathcal{S}\!\mathit{paces}}
\def\Quotfunctor{\mathrm{Quot}}
\def\Hilbfunctor{\mathrm{Hilb}}
\def\Curvesstack{\mathcal{C}\!\mathit{urves}}
\def\Polarizedstack{\mathcal{P}\!\mathit{olarized}}
\def\Complexesstack{\mathcal{C}\!\mathit{omplexes}}
% \Pic is the operator that assigns to X its picard group, usage \Pic(X)
% \Picardstack_{X/B} denotes the Picard stack of X over B
% \Picardfunctor_{X/B} denotes the Picard functor of X over B
\def\Pic{\mathop{\mathrm{Pic}}\nolimits}
\def\Picardstack{\mathcal{P}\!\mathit{ic}}
\def\Picardfunctor{\mathrm{Pic}}
\def\Deformationcategory{\mathcal{D}\!\mathit{ef}}
\usepackage{cleveref}
\crefname{originalsection}{original section}{original sections}
\Crefname{originalsection}{Original section}{Original sections}
\makeatletter
\newlabel{example-graded-vector-spaces}{{0FFX}{source}{}{sourcepointer-1}{}}
\newlabel{example-graded-vector-spaces@cref}{{[originalsection][0][]0FFX}{[1][0][]source}{}{}{}}
\makeatother
\crefname{originaldefinition}{original definition}{original definitions}
\Crefname{originaldefinition}{Original definition}{Original definitions}
\makeatletter
\newlabel{categories-definition-dual}{{0FFP}{source}{}{sourcepointer-2}{}}
\newlabel{categories-definition-dual@cref}{{[originaldefinition][0][]0FFP}{[1][0][]source}{}{}{}}
\makeatother
\title{A graded vector space with a left dual is finite dimensional and the evaluation is a perfect pairing}
\author{The Stacks Project Authors}
\date{stacks-project:homology.tex:local-snapshot-20261012; GNU FDL 1.2 or later}
\begin{document}
\maketitle
\noindent\textit{Editorial note.} All delivered mathematical text is the author's own, copied from the pinned Stacks Project chapter file homology.tex; no statement or proof step was rewritten, reordered or removed.\par\medskip
\noindent\textit{Editorial note.} The proof is the author's complete argument; the short sentence 'As unit we take' that opens it is the author's own text and is preserved verbatim inside the counted span, and the closing sentences derive the statement rather than deferring it.\par\medskip
\noindent\textit{Editorial note.} Every labelled reference occurring inside the counted spans is carried in assembly.external\_references as a non-load-bearing navigation pointer (same-chapter ones at their definition site, cross-chapter ones at the invocation inside the counted span), so no load-bearing mathematics is replaced by a stub. No bibliography citation occurs inside the counted spans.\par\medskip
\noindent\textit{Editorial note.} The same-name check for this card scanned every staged Stacks chapter of this repository -- all sources/stacks-*/pinned-original-source.tex files, the card's own chapter excluded -- and the labels found are recorded in dedup.same\_name\_cross\_chapter on the card.\par\medskip
\noindent\textit{Editorial note.} The preamble is the chapter's own native preamble excerpt already used for this source: the comment-type reference, slogan and history environments (native preamble lines 8-13), the Xy-pic setup (native preamble lines 14-18) and the native theorem declarations with the mathematical macros (native preamble lines 154-217, the actual last line of the file).\par\medskip
\noindent\textit{Allowed standard prerequisites (editorial).}\begin{itemize}
\item The monoidal structure on graded vector spaces. Same chapter, Example example-graded-vector-spaces. Supplies the tensor product, the associativity constraints and the unit of the monoidal category in which the left dual is taken.
\item Dual objects in a monoidal category. Categories, Definition categories-definition-dual. Supplies the evaluation and coevaluation maps together with the two identities that the proof unfolds.
\end{itemize}\medskip
\ifdefined\hypertarget\hypertarget{sourcepointer-1}{}\fi
\noindent\textit{Source reference.} Original section 0FFX: sources/stacks-homology/pinned-original-source.tex, source lines 4251--4251. Navigation pointer invoked at proof line 4284; the counted span contains the pointer only and the cited result is not reproduced inside it, so the pointer is registered as a non-load-bearing navigation stub in the established form for Stacks material. No bibliography citation is involved.\par\medskip
\ifdefined\hypertarget\hypertarget{sourcepointer-2}{}\fi
\noindent\textit{Source reference.} Original definition 0FFP: sources/stacks-categories/pinned-original-source.tex, source lines 9526--9526. Navigation pointer invoked at proof line 4292; the counted span contains the pointer only and the cited result is not reproduced inside it, so the pointer is registered as a non-load-bearing navigation stub in the established form for Stacks material. No bibliography citation is involved.\par\medskip
\begin{lemma}
\label{lemma-left-dual-graded-vector-spaces}
Let $F$ be a field. Let $\mathcal{C}$ be the category of graded
$F$-vector spaces viewed as a monoidal category as in
Example \ref{example-graded-vector-spaces}. If $V$ in $\mathcal{C}$
has a left dual $W$, then $\sum_n \dim_F V^n < \infty$
and the map $\epsilon$ defines nondegenerate pairings
$W^{-n} \times V^n \to F$.
\end{lemma}

\begin{proof}
As unit we take
By Categories, Definition \ref{categories-definition-dual} we have
maps
$$
\eta : \mathbf{1} \to V \otimes W\quad
\epsilon : W \otimes V \to \mathbf{1}
$$
Since $\mathbf{1} = F$ placed in degree $0$, we
may think of $\epsilon$ as a sequence of pairings
$W^{-n} \times V^n \to F$ as in the statement of the lemma.
Choose bases $\{e_{n, i}\}_{i \in I_n}$ for $V^n$ for all $n$.
Write
$$
\eta(1) = \sum e_{n, i} \otimes w_{-n, i}
$$
for some elements $w_{-n, i} \in W^{-n}$ almost all of which are zero!
The condition that $(\epsilon \otimes 1) \circ (1 \otimes \eta)$ is the
identity on $W$ means that
$$
\sum\nolimits_{n, i} \epsilon(w, e_{n, i})w_{-n, i} = w
$$
Thus we see that $W$ is generated as a graded vector space
by the finitely many nonzero vectors $w_{-n, i}$.
The condition that $(1 \otimes \epsilon) \circ (\eta \otimes 1)$
is the identity of $V$ means that
$$
\sum\nolimits_{n, i} e_{n, i}\ \epsilon(w_{-n, i}, v) = v
$$
In particular, setting $v = e_{n, i}$ we conclude that
$\epsilon(w_{-n, i}, e_{n, i'}) = \delta_{ii'}$. Thus
we find that the statement of the lemma holds and that
$\{w_{-n, i}\}_{i \in I_n}$ is the dual basis for $W^{-n}$ to
the chosen basis for $V^n$.
\end{proof}

\end{document}
